\section{RAL-MoE-ROM}
\label{sec:method}

RAL-MoE-ROM separates reduced modeling into three coupled components: regime-local representations, autonomous local specialists, and
cross-regime physical-space assembly. These components are organized by a two-level routing hierarchy. Overlapping regime-local charts provide independent reduced coordinate systems. Each chart supports a complete autonomous physics--data specialist that combines projected ROM dynamics with a structured sparse correction. The outer routing level then selects one specialist or an admissible adjacent pair and assembles only their reconstructed physical fields. Figure~\ref{fig:ral_overview} summarizes the overall construction; detailed
derivations are provided in Appendices~\ref{app:problem_pod}--\ref{app:routing}.

\subsection{Overlapping regime-local reduced charts}

Let $r\in\{S,H,P\}$ index steady, Hopf-onset, and periodic regimes, and let $\mathcal M_r$ denote the corresponding overlapping parameter regions. Each regime is equipped with its own affine reduced coordinate system. A physical velocity field is encoded and reconstructed as
\begin{equation}
  \bm a^{(r)}={\Phi_u^{(r)}}^{\top}M_u(\bm u-\bar{\bm u}^{(r)}),\qquad
  \widehat{\bm u}^{(r)}=\bar{\bm u}^{(r)}+\Phi_u^{(r)}\bm a^{(r)},
  \label{eq:main_chart}
\end{equation}
where $M_u$ is the finite-volume quadrature matrix. Gauge-fixed pressure is encoded analogously, yielding local coefficients
$\bm b^{(r)}$. 

A local chart is therefore defined not only by its basis, but by the full collection of regime-specific quantities used by the reduced model: the parameter region, affine means, velocity and pressure bases, normalization statistics, and projected operators.
All of these quantities are constructed independently across regimes; equal reduced dimensions therefore do not imply componentwise
correspondence between local coordinates. Cross-regime assembly is consequently performed only after reconstruction in physical space.

\subsection{Regime-local physics--data specialists}

Within each chart, projection provides the physics-based dynamical backbone, while a learned closure accounts for unresolved reduced dynamics:
\begin{equation}
  \dot{\bm a}^{(r)}=
  \underbrace{\bm c_r+\bm A_r\bm a^{(r)}+
  \bm H_r:\bigl(\bm a^{(r)}\otimes\bm a^{(r)}\bigr)+
  \bm C^p_r\bm b^{(r)}}_{\mathcal F_r^{G}(\bm a^{(r)},\bm b^{(r)};\mu)}
  +\bm\rho^u_{\theta_r}(\bm\xi^{(r)}).
  \label{eq:main_specialist}
\end{equation}
The learned term $\bm\rho^u_{\theta_r}$ is the velocity closure associated with chart $r$. Its input $\bm\xi^{(r)}$ collects standardized local state variables, the projected right-hand side, the physical parameter, and a short history of preceding states.

Separate sparse MoE maps are used for the velocity correction $\bm\rho^u_{\theta_r}$ and the learned pressure candidate $\bm\rho^p_{\theta_r}$:
\begin{equation}
  \begin{aligned}
  \mathcal M^c_{\theta_r}(\bm\xi)&=
  \omega^c_0 E^c_{0,g^\star}(\bm\xi)+
  \sum_{e\in\mathcal T_2^c}\omega^c_e E^c_{e,g^\star}(\bm\xi),\\
  E(\bm\xi)&=N(\bm\xi)+W_{\rm lin}\bm\xi+
  \kappa W_o\bigl[(U^\top\bm\xi)\odot(V^\top\bm\xi)\bigr].
  \end{aligned}
  \label{eq:main_inner_moe}
\end{equation}
Routing is hierarchical within each specialist. A first router selects one chart-local expert group $g^\star$. Within that group, the velocity and pressure channels $c\in\{u,p\}$ independently select their two highest-scoring routed experts, while a shared expert remains active. The nonnegative weights $\omega_e^c$ are normalized over the active experts. In the expert map, $N$ denotes the nonlinear branch, $\kappa$ scales the bilinear contribution, and $\odot$ denotes componentwise multiplication. The three branches provide nonlinear, affine, and low-rank bilinear response components, while the projected Galerkin term remains explicitly present in the dynamics.

After advancing velocity over one reduced time step, pressure is reconstructed algebraically as
\begin{equation}
  \bm b_{n+1}^{(r)}
  =\bm\gamma_n^{(r)}\odot\mathcal Q_r^{PP}(\bm a_{n+1}^{(r)};\mu)
  +\bigl(\bm 1-\bm\gamma_n^{(r)}\bigr)
   \odot\bm\rho^p_{\theta_r}(\bm\xi_n^{(r)}).
  \label{eq:main_pressure_closure}
\end{equation}
Here $\mathcal Q_r^{PP}$ denotes the chart-local reduced Pressure--Poisson reconstruction, $\bm\rho^p_{\theta_r}$ provides the
learned pressure candidate, and $\bm\gamma_n^{(r)} =\operatorname{sigmoid}(\bm\eta^p_{\theta_r}(\bm\xi_n^{(r)}))$ is a componentwise mixing gate. Thus velocity is advanced dynamically, whereas pressure is reconstructed algebraically rather than evolved as a separate recurrent state.


\subsection{Closed-loop autonomous training}

Each specialist is initialized from an encoded physical history and is then advanced recursively using its own predicted states. To expose the learned correction to the states encountered during autonomous use, training unrolls each specialist for $K$ steps and minimizes the
normalized coefficient error
\begin{equation}
  \mathcal L_r=\frac{1}{K}\sum_{k=1}^{K}\left[
  \frac{\left\|(\widehat{\bm a}_{n+k}^{(r)}-\bm a_{n+k}^{(r)})
  \oslash\bm\sigma_a^{(r)}\right\|_2^2}{d_u}
  +\frac{\left\|(\widehat{\bm b}_{n+k}^{(r)}-\bm b_{n+k}^{(r)})
  \oslash\bm\sigma_b^{(r)}\right\|_2^2}{d_p}\right].
  \label{eq:main_rollout_loss}
\end{equation}
Here $d_u$ and $d_p$ denote the velocity and pressure coordinate dimensions, while $\bm\sigma_a^{(r)}$ and $\bm\sigma_b^{(r)}$ are scaling
vectors estimated from the corresponding regime-local training set; $\oslash$ denotes componentwise division. This coefficient-space objective is used for specialist training, whereas evaluation is performed on reconstructed physical fields. After initialization, no reference state from inside the prediction window is injected into the rollout. The numerical integrator and horizon curriculum are detailed in Appendix~\ref{app:specialists}, and held-out trajectory partitions are reported in Table~\ref{tab:chart_splits}.



\subsection{Cross-regime routing and physical-space T2-C assembly}

The outer E2 router receives only the standardized physical parameter $\widetilde\mu$ and returns a trajectory-level preference over the
regime-local specialists,
\[
\bm\pi(\mu)
=
\operatorname{softmax}\!\left(
\mathcal R_{E2}(\widetilde\mu)
\right).
\]
The probabilities are evaluated once and remain fixed throughout the rollout. The outer router does not use physical history. Chart-independent descriptors extracted from the initial history $\mathcal H_n$ enter only the T2-C correction described below.

A validation-derived binary mask $m_r(\mathcal H_n,\mu,K)$ marks whether specialist $r$ is admissible for the current parameter, initial history, and rollout horizon. The mask restricts the candidate set but does not modify the E2 probability
model itself. Under hard Top-1 routing, only the admissible specialist with the largest E2 probability is advanced.

When an adjacent pair is admissible ($S$--$H$ or $H$--$P$), both specialists are initialized from the same physical history, encoded in
their own coordinates, and advanced independently. Let $\bar\pi_r$ denote the E2 probability normalized over the selected pair. For reconstructed fields $\widehat{\bm q}^{(r)}$ and $\widehat{\bm q}^{(s)}$, T2-C predicts a history-conditioned correction to this pairwise preference and forms
\begin{equation}
  \widehat{\bm q}_{\rm T2C}(t)
  =\alpha\,\widehat{\bm q}^{(r)}(t)
  +(1-\alpha)\,\widehat{\bm q}^{(s)}(t),
  \qquad
  \alpha=
  \sigma\!\left(
  \operatorname{logit}(\bar\pi_r)+\Delta_\psi
  \right).
  \label{eq:main_t2c}
\end{equation}
Here $\Delta_\psi=c_\psi([\widetilde\mu;\bm d_n])$ is a learned correction based on the standardized parameter and the chart-independent physical-history descriptor $\bm d_n$, and $\sigma$ denotes the sigmoid. The scalar $\alpha$ is computed once for the prediction window and is shared by velocity and pressure. The field vector $\bm q$ collects reconstructed velocity and gauge-fixed pressure on the common physical mesh.

The two local trajectories remain independent throughout the rollout: fusion is applied only to their reconstructed outputs and is never fed back into either specialist. If no specialist is admissible, the method abstains from producing a ROM prediction. Admissibility is derived empirically from validation support and therefore does not constitute a formal stability guarantee.
